With two simulated missing sensors, the fine-tuned lifted cavity network improves on FA-Gaussian by
\RairqcoeTwovsfa\ \RairqcoeTwovsfaci\ nats for CO and \RairqnoTwoeTwovsfa\ \RairqnoTwoeTwovsfaci\ for NO$_2$. The
fine-tuned transformer is again better: the lifted network's gain over it is \RairqcoeTwovspfn\ \RairqcoeTwovspfnci\
(CO) and \RairqnoTwoeTwovspfn\ \RairqnoTwoeTwovspfnci\ (NO$_2$). On CO, the cavity input stops helping after
fine-tuning.

\paragraph{Why the transformer wins on real data.}
The lifted network underfits rather than overfits: on the Beijing fine-tuning episodes, the transformer's final
training NLL is \FTbeijingtraingap\ nats lower, about the size of the test gap. Post-hoc variants do not close the gap
on Beijing (Table~\ref{tab:explore}; natural missingness, transformer \RbeijingeZerovalftpfn): a second nuisance dimension scores \Xplifttwo\
and five times more fine-tuning steps \Xpliftlong, against \Xpliftone\ for the same seed with the standard recipe.
Nor is the edge that of a flexible per-episode regression: a Gaussian process fitted to each episode's calibration
rows is the worst method on Air Quality (\RairqcoeZerovalgp\ nats for CO with all sensors observed). These point to
the Gaussian site family itself. Gaussian sites with per-sensor re-linearization are a strong prior for synthetic
factor noise, but real stations share structure they cannot express, and a generic model learns it from a few
thousand real episodes: fine-tuning improves the transformer by \FTbeijingpfngain\ nats on Beijing and the lifted
network by \FTbeijingliftgain.
